% ==============================================================================
% SLIDES - SEMANA 08: PRINCÍPIOS DE OPERAÇÃO DO CLP, CICLO DE SCAN E ENDEREÇAMENTO
% ==============================================================================

% 1. CAPA
\senaicover
  {Princípios de Operação do CLP, Ciclo de SCAN e Endereçamento}
  {Determinismo temporal, ciclo de varredura e mapeamento de memória.}
  {Unidade Curricular: Automação Industrial}
  {Prof. André Souza}

% 2. AGENDA
\begin{senaiframe}{Visão Geral \& Agenda}{Estrutura da Aula}
  \begin{columns}[t]
    \begin{column}{0.48\textwidth}
      \begin{tcolorbox}[senaibluecard, title={1. Conceitos Chave}]
        \begin{itemize}
          \item Fundamentos teóricos e definições.
          \item Relevância para a Indústria 4.0.
        \end{itemize}
      \end{tcolorbox}
      \vspace{4pt}
      \begin{tcolorbox}[senaiambercard, title={3. Aplicação Prática}]
        \begin{itemize}
          \item Estudo de caso na Célula Smart N1.
          \item Desenvolvimento de rotinas técnicas.
        \end{itemize}
      \end{tcolorbox}
    \end{column}
    \begin{column}{0.48\textwidth}
      \begin{tcolorbox}[senairedcard, title={2. Detalhamento Técnico}]
        \begin{itemize}
          \item Ferramentas e protocolos utilizados.
          \item Boas práticas de implementação.
        \end{itemize}
      \end{tcolorbox}
      \vspace{4pt}
      \begin{tcolorbox}[senaigreencard, title={4. Exercícios \& Avaliação}]
        \begin{itemize}
          \item Desafios de fixação do conhecimento.
          \item Integração com a plataforma FactoryHub.
        \end{itemize}
      \end{tcolorbox}
    \end{column}
  \end{columns}
\end{senaiframe}

% 3. CONTEÚDO TEÓRICO
\begin{senaiframe}{Teoria • Módulo 1}{Princípios de Operação do CLP, Ciclo de SCAN e Endereçamento}
  \begin{columns}[t]
    \begin{column}{0.48\textwidth}
      \begin{tcolorbox}[senaibluecard, title={Fundamentos Teóricos}]
        \begin{itemize}
          \item Fases do SCAN (Leitura de Entradas -> Execução da Lógica -> Escrita de Saídas), watchdog timer e endereçamento IEC 61131-3.
        \end{itemize}
      \end{tcolorbox}
    \end{column}
    \begin{column}{0.48\textwidth}
      \begin{tcolorbox}[senaigreencard, title={Prática \& Laboratório}]
        \begin{itemize}
          \item Cálculo do tempo de resposta do ciclo de SCAN.
        \end{itemize}
      \end{tcolorbox}
    \end{column}
  \end{columns}
\end{senaiframe}
